\documentclass[11pt,a4paper]{article}

\usepackage[T1]{fontenc}
\usepackage[utf8]{inputenc}
\usepackage[a4paper,margin=2.8cm]{geometry}
\usepackage{microtype}
\usepackage[hidelinks]{hyperref}

\hypersetup{
  pdftitle={Menelaus}
}

\setlength{\parindent}{1.25em}
\setlength{\parskip}{0.35em}
\raggedbottom

\title{Menelaus}
\author{}
\date{}

\begin{document}

\maketitle

\begin{verse}
I

One of the best moments of the year\\
is when you travel to Sparta.\\
Menelaus organizes a gathering, attended\\
by those who survived Troy and the Nostos.

Of course, you are the guest of honor,\\
and they pamper you with food and wine,\\
and there is a pretty slave with soft hands\\
who takes good care of your tired body.

And you all have a good time, recalling\\
episodes, repeated so often over the years\\
that it is hard to separate truth from fantasy,\\
reality from imagination.

Sometimes you suspect that half your deeds\\
are but reconstructions, as pretty as your slave girl,\\
of something so simple and dirty and brutal,\\
conjured to cover the blood and the dirt and the dung.

But who cares now? You all gather each year,\\
and each year there are fewer and fewer guests at Menelaus's,\\
and eat, and drink, and sing, to celebrate\\
the pretense that your lives made any sense.

II

Fewer and fewer of you. Most are gone.\\
Many of your friends didn't make it back from Troy,\\
Achilles, Ajax, Patroclus, Diomedes,\\
the list is long. Others died on their way back.

Much has been written about the gods\\
opposing the return of the Argives,\\
and little about the rats that, you suspect,\\
were behind the sicknesses that killed so many.

Ah, but people love to sing the gods,\\
nasty and capricious as they are,\\
and to celebrate their meddling in human affairs,\\
their insane lust and cruelty.

Given the choice, you would stick to the rats,\\
which, at least, minded their own business,\\
and although they ate your food and drank your water,\\
they didn't have fun with your ruin, like the Olympians.

III

It has been a long time\\
since you ceased to believe in those people\\
supposedly divine, who drank nectar and had parties\\
on top of Mount Olympus.

Most likely, none of your companions,\\
the few old men still standing,\\
believed in them. But all of you pretend\\
a convenient devotion to keep the status quo.

For, if gods are not real,\\
why should people take kings seriously?\\
Why suffer tyrants like you all,\\
if they are as empty of meaning as the deities?

So you carry on with thigh burnings,\\
and libations, and prayers, dutifully conducted\\
by this other species of rats, the priests,\\
always ready to eat your food and drink your wine.

In exchange, of course, for prophecies\\
that happen to match your own schemes,\\
and signs that confirm your whims,\\
and general blessing of your caprices.

So, it's true, gods do not exist,\\
nor are they necessary, after all,\\
men can take care of abusing,\\
and enslaving, and killing each other.

IV

Those who still meet each year\\
pretend to be brothers of blood,\\
and true, there is warmth between you,\\
something close to friendship.

Ah, but how eagerly\\
each guest gauges the state of the other,\\
measures the ruined skin,\\
the dwindling hair, the broken back.

How precisely you assess\\
the dimmed eyes, the trembling hand,\\
the drooping eyelids, the tremor in the voice,\\
all the signs that scream old age.

This one has grown too fat, this one too thin,\\
Idomeneus sleeps all the time, Philoctetes is always drunk,\\
Nestor can barely speak, most of them\\
are in worse shape than you.

And you wonder if the only meaning of those meetings\\
is to reinvent the past,\\
scare away the encroaching loneliness,\\
and count the surviving bodies.

V

Menelaus is doing pretty well,\\
it can be argued that you and he\\
are the ones least damaged by time,\\
mostly because both are still lucid.

And you actually love him,\\
there is some truth among the many lies,\\
you were brothers in arms,\\
and more than once, you saved each other's lives.

So the bond is real,\\
so real, in fact, that only Telemachus\\
is dearer to your heart,\\
without counting, of course, Circe.

But you won't mention the sorceress\\
to your friend, as he won't mention Helen,\\
both of you have learned\\
to respect each other's broken hearts.

Not that you are teenagers in love,\\
the sorrows of the young are as easy to mend\\
as their strong bodies. But an old man's lost love\\
is the last pain he will carry to the grave.

VI

And so, Menelaus won't speak of Helen,\\
she has been gone for so many years now,\\
in fact she disappeared shortly after\\
the visit of Telemachus to Sparta.

Your friend bragged so much to your son,\\
made sure he saw the fabled Helen in all her splendor,\\
so that he would tell his father Odysseus,\\
in case he had survived.

And indeed, Telemachus said\\
he had never seen anyone more beautiful,\\
and more distant. Looking at her cold majesty,\\
he was convinced of her divine nature.

So when Helen vanished, the official tale\\
spoke of a holy ascent to Elysium,\\
she had not eloped again, no,\\
but had been recalled by the gods.

VII

No one questioned the story,\\
as no one questions your adventure with the Cyclops,\\
tales were invented for this precise purpose,\\
to find sense where there is only noise.

Sound and fury, like the hordes you led,\\
rage, and hate and lust, pumping your hearts,\\
bronze, and blood and fire,\\
to be disguised as epic battles.

So you all went to Troy,\\
to avenge the rape of Menelaus's wife,\\
and fought for ten years, and burned the city,\\
and took the trophy back to the palace.

But the second time she went missing,\\
there was no city to sack, no unity among the Argives,\\
and too many old men in charge,\\
so invoking the gods was appropriate.

VIII

The amazing thing\\
is that Menelaus loved her.\\
Loved her much more after Troy\\
than he ever did before.

After their wedding,\\
Helen was just a trophy, and the prince,\\
so young and furious, bedding other women,\\
was just good, legal sport for a man.

The supposed kidnapping made him mad,\\
perhaps because he knew, like everybody else,\\
that no one had forced Helen to jump\\
into Paris's ship, or into Paris's bed.

Tales tell that in Troy\\
Helen bared her breast to her husband,\\
invited him to strike, and Menelaus\\
dropped his sword and cried.

And for once, you are convinced\\
that this one story is true,\\
Menelaus had needed a refusal and a war\\
to understand he was in love.

IX

But then he made the classic mistake:\\
he assumed that his newly found love,\\
his fervent desire, his admiration for her,\\
would gain her heart.

He was wrong. Helen, of course,\\
would behave as the noblewoman she was,\\
and take care of house and bed,\\
keeping her cards close.

But it was her husband who killed Deiphobus,\\
and you saw her cry her heart out over his body,\\
how ironic, you thought. Not the brave Menelaus,\\
nor the coward Paris, but that other,

discreet, valiant prince,\\
almost as brave and honest as Hector.\\
To him, and only to him,\\
Helen gave her heart.

X

And for him, and only for him,\\
she plotted her revenge.\\
She didn't choose the knife, like her sister,\\
nor the poison all too familiar to her.

Nay, she wanted to hurt more,\\
so she worked hard until Menelaus\\
could only breathe the air she had breathed,\\
and then knifed him in the back.

No bronze, nor even the new iron daggers,\\
was necessary. She simply left,\\
some say that a Trojan called Aeneas\\
organized the rescue party.

XI

Menelaus won't speak of her,\\
as you won't speak of Circe,\\
every year, the two of you gather\\
in Sparta and drink thick wine,

and get drunk, and cry over the lost friend,\\
and the lost love\\
that you never deserved.
\end{verse}

\end{document}
